\documentclass{article}
\usepackage{amsmath,amssymb,amsthm}
\newtheorem{lemma}{Lemma}
\newtheorem{theorem}{Theorem}
\begin{document}
\title{Sharp binary rooftop energy obstruction in original smooth backgrounds}
\date{}\maketitle

The result below concerns a proposed sufficient estimate, not the literal
continuity question itself. It constructs one allowed original smooth
strongly positive background and bounded actual potentials for which
$e(P(u,v))>e(u)+e(v)$. It also proves that the uniform binary constant $2$ is sharp in degree
one across the original background class. This does not exclude constants
depending on a fixed background, direct finite-family bounds, or C1.

\paragraph{Original geometry and frozen inputs.}
Keep the prescribed finite extension $K/\mathbf Q_p$, its unchanged
absolute value, a uniformizer $\pi$, and $\lambda=-\log|\pi|>0$. Set
\[
 X=(\mathbf G_{m,K}^2)^{\rm an},\quad
 t=(\log|T_1|,\log|T_2|),\quad Q=(-\lambda,\lambda)^2,
 \quad V=t^{-1}(Q),
\]
\[
 \beta_t=d't_1\wedge d''t_1+d't_2\wedge d''t_2,
 \qquad S_0=\beta_t\wedge[X].
\]
The closed box lies in the ordinary interior of the larger original
Laurent affinoid $t^{-1}([-2\lambda,2\lambda]^2)$. All backgrounds below
extend to its ordinary neighborhood. We use the unchanged relative
strong completion with its encoded approximation boundary, and retain
the full original test sheaf. Here $n=2$, $p=q=1$ and $d=2$.
Frozen R17 FP1--FP6 supplies canonical finite convex Gauss traces,
their completed linear moment identity, and compact toric test pairings
for $S_0$. IF20--IF21 supplies full original toric calibration and
quadratic energy equality for bounded Lipschitz convex zero-boundary
potentials. The actual-core admission argument following IF20 supplies
finite mass, compact negative sublevels and canonical no loss for these
neighborhood-extendible functions. Frozen HC14--HC15 places every original
compact smooth scalar in the actual difference space. Finite relative
rooftops exist and have their frozen greatestness property.

\begin{lemma}[Identification of a bounded polyhedral toric rooftop]
Let $f_1,f_2$ be bounded Lipschitz convex functions on
$Q_0=(-1,1)^2$, continuous and zero on its boundary, negative inside,
with convex neighborhood extensions. Suppose their real convex envelope
$h=\operatorname{conv}(\min(f_1,f_2))$ has the same properties. Set
$F_i=f_i(t/\lambda)$ and $H=h(t/\lambda)$. Then the completed relative
rooftop for $S_0$ equals $H$ in the original relative quotient.
\end{lemma}
\begin{proof}
The IF20 actual-core argument admits all three functions. Their bounded
convex neighborhood extensions have finite Laplacian mass, compact strict
negative sublevels, finite energies and eventually unchanged deep
truncations. Thus $H\preceq P:=P_{S_0}(F_1,F_2)\preceq F_i$.
To justify trace order without evaluating a relative class at exceptional
Gauss points, replace any Borel representative of $P$ by its pointwise
clipping between $H$ and $\min(F_1,F_2)$. This agrees with $P$ outside
the fixed union of allowed relative null sets and hence represents the
same completed member. It satisfies the two order bounds everywhere.
The canonical convex trace $g_P$ agrees with this representative on the
Gauss section almost everywhere. Therefore
\[
 h(z)\le g_P(\lambda z)\le\min(f_1(z),f_2(z))
\]
almost everywhere, and hence everywhere by convex continuity. Greatestness
of the real convex envelope forces $g_P(\lambda z)=h(z)$.

This trace equality alone is not used as an injectivity assertion.
For every compact smooth real $\varphi$ on $Q$, frozen FP6 gives
\[
 B_{S_0}(P-H,\varphi\circ t)
       =-\int_Q(g_P-h(t/\lambda))\Delta\varphi\,dt=0.       \tag{1}
\]
Choose compact smooth $\varphi_k$ converging to $h(t/\lambda)$ in
$H_0^1(Q)$. Such approximation applies because $h$ is Lipschitz and
zero on the rectangle boundary. The compact pullbacks are actual energy
differences. Expanding their difference energy by FP6 and IF21 gives
\[
 B_{S_0}(\varphi_k\circ t-H,\varphi_k\circ t-H)
       =\int_Q|\nabla\varphi_k-\nabla(h(t/\lambda))|^2dt
          \longrightarrow0.                              \tag{2}
\]
Quadratic energy Cauchy--Schwarz and (1) imply $B_{S_0}(P-H,H)=0$.
Consequently
\[
 e_{S_0}(P)=e_{S_0}(H)+I_{S_0}(P,H).
\]
But $H\preceq P$ gives $e_{S_0}(P)\le e_{S_0}(H)$ by frozen energy
monotonicity. Thus $I_{S_0}(P,H)=0$ and the completed potentials are
the same relative member. No identification from trace equality alone,
and no restriction of the original test family, has occurred.
\end{proof}

\begin{lemma}[Uniformly comparable smooth backgrounds]
Let $S=A(t)\wedge[X]$ be smooth, separately closed and positive on a
neighborhood of the closed box, with constants $0<c\le C<\infty$ such that
\[
                    c\beta_t\le A\le C\beta_t.             \tag{3}
\]
Then its degree-one strong cone, relative null ideal and relative order
identify canonically with those for $S_0$. In particular their relative
rooftops agree.
\end{lemma}
\begin{proof}
The degree-one source of each fixed normalized tester is linear in the
background. Applying (3) to every tester, then taking suprema and ordinary-open
outer regularizations, gives two-sided capacity comparisons and hence
equality of the null ideals. The same source
comparisons for bounded actual psh functions give equivalence of finite
mass and power-energy bounds. Compact negative sublevels and the ordinary
boundary conditions are unchanged. The canonical no-loss estimates for
the truncations transfer by the same fixed constants, so the bounded
actual defining cores coincide with their printed qualifications.

For every difference $a-b$ of members of this common actual core, global
Green exchange and positivity of the quadratic energy forms for
$S-cS_0$ and $CS_0-S$ give
\[
 c I_{S_0}(a,b)\le I_S(a,b)\le C I_{S_0}(a,b).              \tag{4}
\]
Both residual backgrounds remain original positive separately closed
backgrounds. Thus the strong Cauchy defining rows and their equivalences
are the same. The capacity comparisons identify their limiting relative
Borel realizations. This proves equality of the completed cones with
their encoded approximation boundary, rather than imposing an ordinary
boundary value on an arbitrary completed member. Their relative orders
and admissible lower families are consequently identical, so greatestness
identifies their rooftops.
\end{proof}

\begin{theorem}
In the original degree-one scope $q=1$, the universal binary estimate
\[
                       e(P(U,V))\le2\bigl(e(U)+e(V)\bigr)
\]
has optimal constant $2$, even when the backgrounds are smooth strongly
positive and the two obstacles belong to the bounded actual core.
Precisely, for every $0<C<2$ there are one allowed original smooth
background and bounded actual obstacles with
\[
                       e(P(U,V))>C\bigl(e(U)+e(V)\bigr).
\]
The background is allowed to depend on $C$. No assertion of sharpness
for a single fixed smooth background is made.
\end{theorem}

\begin{proof}
The upper bound follows from the frozen rooftop root inequality
$r(P(U,V))\le r(U)+r(V)$ and
$(r(U)+r(V))^2\le2(r(U)^2+r(V)^2)$, since $d=q+1=2$.

For sharpness fix $0<a<1$, and work in normalized coordinates
$(x,y)=t/\lambda\in Q_0=(-1,1)^2$. Put
\[
 u_a(x,y)=\max\left\{
 \frac{x-1}{1-a},-\frac{x+1}{1+a},
 \frac{y-1}{1+a},-\frac{y+1}{1-a}\right\},
 \qquad v_a(x,y)=u_a(y,x).
\]
These are convex pyramids over the four zero-height corners, with apices
$(a,-a,-1)$ and $(-a,a,-1)$, respectively. Their common real convex
envelope is
\[
 h_a(x,y)=\max\left\{
 \frac{|x|-1}{1-a},\frac{|y|-1}{1-a},\frac{|x+y|}{2}-1
                         \right\}.                         \tag{5}
\]
To verify this, denote the two planar apices by $A=(a,-a)$ and
$B=(-a,a)$ and the corners by $C_{\pm\pm}=(\pm1,\pm1)$.
The six affine branches in (5) are lower supporting planes through
the triples
\[
 (C_{--},C_{-+},B),\ (C_{--},C_{+-},A),\
 (C_{-+},C_{++},B),\ (C_{+-},C_{++},A),
\]
\[
                   (C_{--},A,B),\ (C_{++},A,B).
\]
They are below all six vertex heights and their projected triangles
tile the square. Affine interpolation on each triangle proves both
the common-minorant inequality and greatestness, by the same elementary supporting-plane argument. All three functions are
globally polyhedral convex, negative inside the square, Lipschitz,
continuous and zero on its boundary. Their original moment pullbacks
are therefore qualified bounded actual core potentials.

On the diagonal one has the exact profiles
\[
 u_a(s,s)=v_a(s,s)=\frac{|s|-1}{1+a},\qquad
 h_a(s,s)=|s|-1.
\]
Consequently, for $E_\ell(f)=\int_{-1}^1|(d/ds)f(s,s)|^2ds$,
\[
 E_\ell(u_a)=E_\ell(v_a)=\frac2{(1+a)^2},\qquad
 E_\ell(h_a)=2,
\]
and
\[
 \frac{E_\ell(h_a)}{E_\ell(u_a)+E_\ell(v_a)}
                         =\frac{(1+a)^2}{2}\longrightarrow2
                         \quad(a\uparrow1).                \tag{6}
\]

Use the same entirely smooth original backgrounds
\[
 S_{\delta,\varepsilon}
 =\left[\varepsilon\beta_t+
 \lambda^2D\!\left(\sqrt{((t_2-t_1)/\lambda)^2+\delta^2}\right)
                                                   \right]\wedge[X],
 \qquad\delta,\varepsilon>0.
\]
They satisfy
$\varepsilon S_0\le S_{\delta,\varepsilon}
\le(\varepsilon+2/\delta)S_0$ on the original neighborhood.
The preceding lemmas identify the completed rooftop of the two actual
obstacles with $h_a(t/\lambda)$ for every positive pair of parameters.
For $f\in\{u_a,v_a,h_a\}$, full original torus calibration and
global Green exchange give the following exact smooth energy formula.
Indeed the background matrix in the $t$ coordinates is
$\varepsilon I+\chi_\delta''(y-x)
\left(\begin{smallmatrix}1&-1\\-1&1\end{smallmatrix}\right)$;
its cofactor is
$\varepsilon I+\chi_\delta''(y-x)
\left(\begin{smallmatrix}1&1\\1&1\end{smallmatrix}\right)$.
The factor $\lambda^2$ in the background definition makes the change
$t=\lambda(x,y)$ cancel in the energy integral. Local convex mollification
as in IF20 passes the calibration on all original tests to each
polyhedral potential; the zero-boundary compact clipping argument of
IF21 justifies global Green exchange. All coefficients are bounded and
smooth for fixed positive parameters. Thus
\[
 e_{\delta,\varepsilon}(f(t/\lambda))
 =\varepsilon\int_{Q_0}|\nabla f|^2dxdy
 +\int_{Q_0}\frac{\delta^2(f_x+f_y)^2}
                       {((y-x)^2+\delta^2)^{3/2}}\,dxdy.    \tag{7}
\]
For fixed $a$ the gradients are bounded. With $r=y-x$, the inner
cross-sectional integral of $(f_x+f_y)^2$ is continuous at $r=0$:
away from finitely many diagonal breakpoints the active plane is locally
fixed; if a whole interface lies on the diagonal, its adjacent affine
functions have the same tangential derivative. Dominated convergence
therefore applies to the cross sections. The nonnegative kernel in (7)
has total real-line mass $2$ and concentrates at zero, so
\[
 \lim_{\delta\downarrow0}e_{\delta,\varepsilon}(f(t/\lambda))
              =\varepsilon E_0(f)+2E_\ell(f),\qquad
 E_0(f)=\int_{Q_0}|\nabla f|^2dxdy.                         \tag{8}
\]
The limiting denominator below is strictly positive, and (6)--(8) give
\[
 \lim_{\varepsilon\downarrow0}\lim_{\delta\downarrow0}
 \frac{e_{\delta,\varepsilon}(h_a(t/\lambda))}
 {e_{\delta,\varepsilon}(u_a(t/\lambda))+
                         e_{\delta,\varepsilon}(v_a(t/\lambda))}
                         =\frac{(1+a)^2}{2}.                \tag{9}
\]
Given $C<2$, choose $a<1$ with the right side greater than $C$,
then choose a sufficiently small positive $\varepsilon$, followed by
a sufficiently small positive $\delta$. Fix these parameters. The
resulting single original smooth background and two actual obstacles
give the required strict inequality. No potential or completion is
transferred to a singular limiting background.
\end{proof}


\paragraph{An explicit fixed smooth counterexample.}
The family above contains a particularly simple fixed example. Take
$a=1/2$ and $\varepsilon=\delta=1/100$. Then
\[
 u=\max\{2(x-1),-\tfrac23(x+1),\tfrac23(y-1),-2(y+1)\},
 \qquad v(x,y)=u(y,x),
\]
\[
 h=\max\{2(|x|-1),2(|y|-1),\tfrac12|x+y|-1\}.
\]
For the single original smooth background $S=S_{1/100,1/100}$ defined
above, the completed rooftop
of $U=u(t/\lambda)$ and $V_1=v(t/\lambda)$ is $H=h(t/\lambda)$.
The background form is bounded between $(1/100)\beta_t$ and
$(200+1/100)\beta_t$. Its full original smooth energy formula is (7).

Here is an exact evaluation rather than a numerical or limiting check.
For a triangle $T$ of area $A_T$ whose three vertex values of $L=y-x$
are $a_T<b_T<c_T$, affine cross-section integration gives
\[
 \int_T\chi_\delta''(L)
 =\frac{2A_T}{c_T-a_T}
  \left(\frac{\chi_\delta(c_T)-\chi_\delta(b_T)}{c_T-b_T}
       -\frac{\chi_\delta(b_T)-\chi_\delta(a_T)}{b_T-a_T}\right),
 \qquad \chi_\delta(r)=\sqrt{r^2+\delta^2}.                 \tag{10}
\]
For a plane $f=\alpha x+\gamma y+\kappa$ on $T$, its energy contribution
is $\varepsilon A_T(\alpha^2+\gamma^2)$ plus $(\alpha+\gamma)^2$
times (10). The four triangles joining each apex to its four boundary
edges give the obstacle energies; the six triangles in the theorem's
proof give the rooftop energy. All relevant vertex values of $L$ are
distinct. Exact substitution gives
\begin{align*}
 e_S(U)=e_S(V_1)
     &=-\frac{16\sqrt{10001}}{225}+\frac2{25}
                                  +\frac{2\sqrt{40001}}{45},\\
 e_S(H)&=-\frac{3\sqrt{10001}}{25}+\frac{13}{100}
                                  +\frac{2\sqrt{40001}}{25}.
\end{align*}
Consequently
\[
 e_S(H)-e_S(U)-e_S(V_1)
     =\frac{\sqrt{10001}}{45}
         -\frac{2\sqrt{40001}}{225}-\frac3{100}
     >\frac{100}{45}-\frac{2\cdot201}{225}-\frac3{100}
      =\frac{73}{180}>\frac25.                            \tag{11}
\]
The strict inequalities use only $\sqrt{10001}>100$ and
$\sqrt{40001}<201$. Thus coefficient-one energy subadditivity fails in
this one explicit smooth original background, with an exact positive gap.

\paragraph{Scope.}
The sharpness statement excludes a background-independent universal binary
constant strictly below $2$ across the original degree-one class. It does
not exclude a subcritical constant $C(S)<2$ for each fixed smooth
background, or a direct arbitrary-finite-family estimate
$e(P(U_i:i\in F))\le C(S,n)\sum_{i\in F}e(U_i)$ whose constant is
independent of $|F|$. It does not disprove SR8 or produce a strongly
convergent row with literal discontinuity. C1 remains open in the full
arbitrary-background scope on the basis of this argument.

\paragraph{Exact frozen source identity.}
The R17 whole source for FP1--FP6, IF20--IF21, and the actual-core and
compact-test qualifications has SHA-256
\texttt{8389d2ce81aa5fd1b563767a3c207a06d0b50e661fb06ee978da2d39634c49be}.
The finite relative rooftop definition and greatestness are recorded
in frozen R12 LG2, SHA-256
\texttt{5001a36a41048e2c6c94189dffcec2f0610537d90ddb64e8c0932770c861988d}.
The completion and positive quadratic energy form remain frozen inputs
in their exact original scopes. The explicit smooth energies above were
independently recomputed from the printed triangulations with exact
arithmetic; the strict final gap uses the displayed rational bounds.
\end{document}
